% titel: Blatt 7, A8: NH3 - ideale Gasgleichung vs. Van-der-Waals-Gleichung
% skript: Mortimer Kap. 11 (reales Gas, Van-der-Waals-Gleichung)

% Teil I: Werkzeuge, kein L\"osungsweg.

Das ideale Gasgesetz vernachl\"assigt das Eigenvolumen der Teilchen und
die gegenseitige Anziehung. Die Van-der-Waals-Gleichung korrigiert beide
Effekte mit den Stoffkonstanten $a$ und $b$.

\begin{werkzeugkasten}[Van-der-Waals-Gleichung]
\[
  \left(p + a\frac{n^2}{V^2}\right)(V - nb) = nRT,
  \qquad
  p = \frac{nRT}{V-nb} - a\frac{n^2}{V^2}.
\]
\begin{itemize}
  \item $a\,(n/V)^2$: Korrektur f\"ur intermolekulare Anziehung
        (zieht $\rightarrow$ senkt $p$).
  \item $b$: Korrektur f\"ur das Eigenvolumen der Molek\"ule
        ($V-nb$ kleiner $\rightarrow$ erh\"oht $p$).
\end{itemize}
\textit{Skript:} Mortimer Kap. 11, reale Gase.
\end{werkzeugkasten}

\begin{achtung}[Darauf achten]
\begin{itemize}
  \item In $p=p(nRT)/(V-nb)-a n^2/V^2$ steht $n^2$ (nicht $n$) beim
        Anziehungsterm.
  \item Einheiten: $a$ in $\mathrm{L^2\,atm\,mol^{-2}}$ und $b$ in
        $\mathrm{L\,mol^{-1}}$ passend zu $p$ in atm und $V$ in L.
  \item Ideal und Van-der-Waals weichen bei dichtem Gas spürbar ab; sind
        die Werte nah beieinander, ist das Gas (bei diesen Bedingungen)
        nahezu ideal.
\end{itemize}
\end{achtung}
